% !TeX program = pdflatex
% Counterexamples to a conjectured characterization of cubes by the lattice
% Borsuk number, formally verified in Lean 4.
% Build: pdflatex note.tex (run twice). See BUILD.md.
\documentclass[11pt]{article}

\usepackage[T1]{fontenc}
\usepackage{lmodern}
\usepackage[margin=1.1in]{geometry}
\usepackage{amsmath,amssymb,amsthm}
\usepackage{booktabs}
\usepackage{microtype}
\usepackage{listings}
\usepackage[hidelinks]{hyperref}
\usepackage{xurl}

\newtheorem{theorem}{Theorem}[section]
\newtheorem{lemma}[theorem]{Lemma}
\newtheorem{conjecture}[theorem]{Conjecture}
\theoremstyle{remark}
\newtheorem{remark}[theorem]{Remark}

\newcommand{\Z}{\mathbb{Z}}
\newcommand{\R}{\mathbb{R}}
\newcommand{\N}{\mathbb{N}}
\newcommand{\conv}{\operatorname{conv}}
\newcommand{\vol}{\operatorname{vol}}
\newcommand{\diamZ}{\operatorname{diam}_{\Z}}
\newcommand{\betaZ}{\beta_{\Z}}
\newcommand{\GL}{\operatorname{GL}}

\lstdefinelanguage{lean4}{
  keywords={theorem,def,abbrev,structure,namespace,end,import,Prop},
  sensitive=true,
  comment=[l]{--},
}
\lstset{
  language=lean4,
  basicstyle=\ttfamily\small,
  keywordstyle=\bfseries,
  columns=fullflexible,
  keepspaces=true,
  showstringspaces=false,
  literate={∧}{{$\wedge$}}1 {∀}{{$\forall$}}1 {ℕ}{{$\mathbb{N}$}}1
           {¬}{{$\neg$}}1 {ℝ}{{$\mathbb{R}$}}1 {∃}{{$\exists$}}1
           {↔}{{$\leftrightarrow$}}1 {ℤ}{{$\mathbb{Z}$}}1
           {×}{{$\times$}}1 {β}{{$\beta$}}1,
}

\title{Counterexamples to a conjectured characterization of cubes\\
by the lattice Borsuk number, formally verified in Lean~4}
\author{John Erlbacher\thanks{Independent Researcher. Email: \texttt{erlbacher.research@gmail.com}. ORCID: 0009-0003-6851-4139.}}
\date{June 2026}

\begin{document}
\maketitle

\begin{abstract}
For a bounded set $S \subset \Z^d$, the lattice Borsuk number $\betaZ(S)$ is
the least number of parts in a partition of $S$ into parts of strictly
smaller lattice diameter. Brose, De Loera, Lopez-Campos, and Torres proved
$\betaZ(S) \le 2^d$ and conjectured (arXiv:2508.20009, Conjecture~3) that
$\betaZ(S) = 2^d$ if and only if $\conv(S)$ is unimodularly equivalent to a
cube $[0,m]^d$. We show that the conjecture is false as stated: the
four-point set $S_A = \{(0,0),(1,0),(0,1),(3,5)\} \subset \Z^2$ has all six
pairwise differences primitive, hence $\betaZ(S_A) = 4 = 2^2$, while
$\conv(S_A)$ contains exactly $7$ lattice points and so is unimodularly
equivalent to no square: the count is a unimodular invariant and $7$ is not a
perfect square. This disproof is formally verified: a Lean~4 development of
about $1100$ lines over mathlib proves \texttt{betaZ SA = 4},
\texttt{latticeCount hullSA = 7}, and \texttt{conjecture3\_false}, with the
formal definitions documented against the verbatim text of the paper; the
axiom audit reports only the three standard axioms. Two further
counterexamples, verified in exact arithmetic but not formalized, close the
natural repair: a four-point set $T = \conv(T) \cap \Z^2$ whose hull is a
triangle of area $3/2$, and an eight-point set $C = \conv(C) \cap \Z^3$ with
$\betaZ(C) = 8 = 2^3$ whose hull has volume $5/2$. The conjecture therefore
fails even when restricted to full sets of lattice points of lattice
polytopes, in dimensions $2$ and~$3$.
\end{abstract}

\section{Introduction}\label{sec:intro}

In a recent study of lattice diameters, Brose, De Loera, Lopez-Campos, and
Torres \cite{BDLT2025} introduced a discrete analogue of Borsuk's classical
partition problem \cite{Borsuk1933,KahnKalai1993}. Throughout, $d \ge 1$ and
all notions are those of \cite{BDLT2025}. For $x, y \in \Z^d$ write
$[x,y] := \conv(\{x,y\})$ and
\[
f(x,y) \;=\; \bigl|[x,y] \cap \Z^d\bigr| - 1 ,
\]
the number of lattice points on the segment, minus one. For a bounded
(equivalently, finite) set $S \subset \Z^d$ the \emph{lattice diameter} is
$\diamZ(S) = \max_{x,y \in S} f(x,y)$, and the \emph{lattice Borsuk number}
$\betaZ(S)$ is the minimal size of a partition of $S$ in which every part has
lattice diameter strictly smaller than $\diamZ(S)$ \cite[Definition~1.3 and
\S 4.2]{BDLT2025}. A vector $u \in \Z^d$ is \emph{primitive} if the gcd of
its coordinates is $1$; two sets $P, Q \subseteq \R^d$ are \emph{unimodularly
equivalent} if $P = AQ + t$ for some $A \in \GL(d,\Z)$ and $t \in \Z^d$ (the
paper states this notion for lattice polytopes; every set to which we apply
it below is one --- see Section~\ref{sec:lean}).

Theorem~4.8 of \cite{BDLT2025} solves the discrete Borsuk partition problem:
$\betaZ(S) \le 2^d$ for every bounded $S \subset \Z^d$, with equality for
$S = \{0,1\}^d$; the proof bounds the maximum degree of the \emph{lattice
Borsuk graph} (vertex set $S$, edges the pairs realizing $\diamZ(S)$) by
$2^d - 1$ via a collinearity lemma of Rabinowitz \cite{Rabinowitz1989} and
applies Brooks' theorem, using that $\betaZ(S)$ is the chromatic number of
this graph. The paper closes with a conjectured characterization of the
equality case, which we quote verbatim (Conjecture~3 of \cite{BDLT2025},
\S 4.2):

\begin{quote}
\emph{``Let $S \subset \Z^d$ be a bounded set. Then $\betaZ(S) = 2^d$ if and
only if $\conv(S)$ is unimodularly equivalent to a $d$-cube $[0,m]^d$ for any
$m \in \N$.''}
\end{quote}

\noindent
(The phrase ``for any $m \in \N$'' is read as ``for some $m \in \N$''; our
counterexamples refute the equivalence under either reading, and for $\N$
with or without $0$ --- see Section~\ref{sec:lean}.) This note proves:

\begin{theorem}[formally verified; see Theorem~\ref{thm:witnessA} and
Section~\ref{sec:lean}]\label{thm:main}
Conjecture~3 of \cite{BDLT2025} is false as stated: its $d=2$ instance fails.
The set $S_A = \{(0,0),\,(1,0),\,(0,1),\,(3,5)\} \subset \Z^2$ satisfies
$\betaZ(S_A) = 4 = 2^2$, yet $\conv(S_A)$ is unimodularly equivalent to no
square $[0,m]^2$, $m \in \N$.
\end{theorem}

\begin{theorem}[exact arithmetic, not formalized; see
Theorems~\ref{thm:witnessB} and \ref{thm:witnessC}]\label{thm:repair}
The conjecture fails even under the additional hypothesis
$S = \conv(S) \cap \Z^d$, in dimensions $2$ and $3$: the sets
\begin{align*}
T &= \{(0,1),\,(1,0),\,(1,1),\,(2,2)\} \subset \Z^2 \quad\text{and}\\
C &= \{(0,2,2),\,(1,0,1),\,(1,1,2),\,(1,2,2),\,(2,1,1),\,(2,1,2),\,(2,2,1),
\,(3,1,3)\} \subset \Z^3
\end{align*}
are the full lattice-point sets of their hulls, attain
$\betaZ = 2^d$, and their hulls are unimodularly equivalent to no cube.
\end{theorem}

In all three witnesses it is the ``only if'' direction of the conjecture that
fails; the upper bound $\betaZ(S) \leq 2^d$ of \cite[Theorem~4.8]{BDLT2025}
is of course untouched. The mechanism is the degenerate regime
$\diamZ(S) = 1$: there every part of an admissible partition must be a
singleton, so $\betaZ(S) = |S|$, and $\betaZ(S) = 2^d$ merely says that $S$
consists of $2^d$ points with all pairwise differences primitive --- a
condition that places almost no constraint on the shape of $\conv(S)$. The
sharpness example $\{0,1\}^d$ of \cite[Theorem~4.8]{BDLT2025} itself lives in
this regime, and the lattice Borsuk graph of each of our witnesses is the
complete graph $K_{2^d}$, exactly the situation the conjecture was designed
to characterize; the failure is therefore not an artifact of a perverse
reading. Theorem~\ref{thm:repair} matters because the authors remark
\cite[\S 4.2]{BDLT2025} that in dimension two the conjecture ``is true and
can be proven using a characterization of lattice polygons which have four
diameter directions'' of B\'ar\'any and F\"uredi \cite[Theorem~2(iii)]
{BaranyFuredi2001} --- a statement about lattice polygons, i.e.\ about sets
of the form $\conv(S) \cap \Z^2$. Witness $T$ shows that no reading of
``bounded set'' as ``full lattice-point set of a polygon'' can save the
$d = 2$ claim: any correct repair must also exclude the diameter-one regime.

What is and is not machine-checked deserves precision.
Theorem~\ref{thm:main} is proved in Lean~4 \cite{Lean4} over mathlib
\cite{mathlib}: the statement of the conjecture is formalized in
\texttt{Borsuk/Defs.lean} with every definition documented against the
verbatim text of \cite{BDLT2025}, and the refuting theorem
\texttt{conjecture3\_false} is checked by the Lean kernel, with an axiom
audit reporting only \texttt{propext}, \texttt{Classical.choice}, and
\texttt{Quot.sound} (Section~\ref{sec:lean}). Theorem~\ref{thm:repair} is \emph{not} formalized;
it is verified by exact integer and rational arithmetic (no floating point)
in a short self-contained script distributed with the artifact and re-run
while preparing this note. A literature check (June 2026, most recently on
2026-06-12) found no prior or concurrent resolution of the conjecture;
arXiv:2508.20009 remains at version~1.

\paragraph{Methods and AI-use disclosure.}
The counterexamples and the formal proof were produced by Demonstrandum, a
verification-first multi-agent AI workflow built on Claude language models
(Anthropic). Agents froze the
conjecture statement from the arXiv v1 \LaTeX{} source, constructed witness
$A$ and three independent $d=3$ full-set witnesses (of which $C$ is one) by
direct search with exact verification, and wrote the Lean development in four
work packages against frozen formal definitions. An agent from a different
model family (Codex, GPT-5.5, OpenAI) independently re-verified all witness
computations, contributed witness $T$, searched independently for prior
resolutions, and adversarially reviewed the faithfulness of the formal
definitions against the paper text. The Lean kernel --- not any language
model --- is the guarantor of every formalized claim; faithfulness of the
formal \emph{definitions} to the paper's statement is the one link the kernel
cannot check, and Section~\ref{sec:lean} documents how it was controlled.
The text of this note was likewise drafted with the Claude-based pipeline;
the author directed the work and takes responsibility for its content. Every
number appearing in this note was recomputed during its preparation (June
2026): the Lean project was rebuilt and its axiom audit re-run, and the
witness data were re-verified by the exact-arithmetic scripts
\texttt{\_recompute.py} and \texttt{\_referee\_recheck.py} included with the
artifact (Section~\ref{sec:verification}).

\section{Witness A and the disproof}\label{sec:witnessA}

The two lemmas below are stated for general $d$; the Lean formalization
instantiates them at $d = 2$ (Section~\ref{sec:lean}), which suffices, since
Conjecture~3 asserts all dimensions at once and witness $C$ is handled
separately. For $v \in \Z^d$ let $\gcd(v)$ denote the gcd of the coordinates
of $v$, with $\gcd(0) = 0$.

\begin{lemma}[segment count; \texttt{Borsuk.ncard\_segLatticePts}]
\label{lem:gcd}
For $x, y \in \Z^d$ with $x \neq y$, let $g = \gcd(y - x) \ge 1$ and
$u = (y-x)/g$. Then
\[
[x,y] \cap \Z^d \;=\; \{\, x + k u : k = 0, 1, \dots, g \,\},
\]
so $\bigl|[x,y] \cap \Z^d\bigr| = g + 1$; for $x = y$ the segment is the
single lattice point $x$ and $\gcd(0) = 0$, so in all cases
$f(x,y) = \gcd(y - x)$. In particular $f(x,y) = 1$ exactly when $y - x$ is
primitive.
\end{lemma}

\begin{proof}
The vector $u$ is primitive and
$x + ku = \bigl(1 - \tfrac{k}{g}\bigr) x + \tfrac{k}{g}\, y \in [x,y]$ for
$0 \le k \le g$, giving $\supseteq$. Conversely, a point $p \in [x,y]$ has
the form $p = x + t(y - x) = x + (tg)u$ with $t \in [0,1]$. By B\'ezout there
are integers $\alpha_1, \dots, \alpha_d$ with $\sum_i \alpha_i u_i = 1$;
applying them to the integer vector $p - x = (tg)u$ gives
$tg = \sum_i \alpha_i (tg) u_i \in \Z$, and $0 \le tg \le g$, so
$p = x + ku$ with $k = tg \in \{0, \dots, g\}$. The map $k \mapsto x + ku$ is
injective ($u \ne 0$), so the count is $g + 1$.
\end{proof}

\begin{lemma}[diameter-one sets;
\texttt{Borsuk.betaZ\_eq\_card\_of\_latticeDiam\_eq\_one}]\label{lem:beta}
Let $S \subset \Z^d$ be finite with $|S| \ge 2$ and all pairwise differences
of distinct points primitive. Then $\diamZ(S) = 1$ and $\betaZ(S) = |S|$.
\end{lemma}

\begin{proof}
By Lemma~\ref{lem:gcd}, $f(x,y) = 1$ for all distinct $x, y \in S$, so
$\diamZ(S) = 1$. In a partition admissible for $\betaZ$, every part has
lattice diameter $< 1$, i.e.\ $0$; but a part containing distinct points
$x \ne y$ has diameter at least $f(x,y) \ge 1$, since distinct lattice points
always have $\gcd(y - x) \geq 1$. Hence all parts are singletons and there
are at least $|S|$ of them; the partition into singletons is admissible
($0 < 1$), so $\betaZ(S) = |S|$.
\end{proof}

\begin{theorem}[formalized: \texttt{witnessA\_kills\_conjecture3},
\texttt{conjecture3\_false}]\label{thm:witnessA}
Let $S_A = \{(0,0),\,(1,0),\,(0,1),\,(3,5)\} \subset \Z^2$. Then:
\begin{enumerate}\itemsep2pt
\item[(i)] $\betaZ(S_A) = 4 = 2^2$;
\item[(ii)] $\conv(S_A) \cap \Z^2$ consists of exactly the $7$ points
$(0,0)$, $(1,0)$, $(0,1)$, $(1,1)$, $(1,2)$, $(2,3)$, $(3,5)$;
\item[(iii)] for every $m \in \N$, $\conv(S_A)$ is not unimodularly
equivalent to $[0,m]^2$.
\end{enumerate}
In particular the $d = 2$ instance of Conjecture~3 of \cite{BDLT2025} ---
hence the conjecture, which asserts all dimensions --- is false.
\end{theorem}

\begin{proof}
(i) The six pairwise differences are $(1,0)$, $(0,1)$, $(3,5)$, $(-1,1)$,
$(2,5)$, $(3,4)$, each with coprime coordinates; apply
Lemma~\ref{lem:beta}.

(ii) Let
$Q = \{(x,y) \in \R^2 : x \ge 0,\ y \ge 0,\ 5x - 2y \le 5,\ 3y - 4x \le 3\}$.
Each of the four halfplanes is convex and contains the four points of $S_A$
(sixteen evaluations; each bounding line passes through exactly two points of
$S_A$, namely the edge pairs $\{(0,1),(0,0)\}$, $\{(0,0),(1,0)\}$,
$\{(1,0),(3,5)\}$, $\{(0,1),(3,5)\}$), so $\conv(S_A) \subseteq Q$. For
integer points of $Q$: the third and fourth inequalities give
$15x - 15 \le 6y \le 8x + 6$, so $x \le 3$, and then $y$ is pinned down case
by case --- $x = 0$ forces $y \in \{0,1\}$; $x = 1$ forces $y \in \{0,1,2\}$;
$x = 2$ forces $y = 3$; $x = 3$ forces $y = 5$. These are the seven listed
points. Conversely all seven lie in $\conv(S_A)$: four are points of $S_A$,
and
\begin{gather*}
(1,1) = \tfrac{2}{5}(0,0) + \tfrac{2}{5}(1,0) + \tfrac{1}{5}(3,5), \qquad
(1,2) = \tfrac{1}{3}\bigl[(0,0) + (0,1) + (3,5)\bigr],\\
(2,3) = \tfrac{1}{5}(0,0) + \tfrac{1}{5}(1,0) + \tfrac{3}{5}(3,5).
\end{gather*}

(iii) If $\Phi(x) = Ax + t$ with $A \in \GL(2,\Z)$, $t \in \Z^2$, then $\Phi$
and $\Phi^{-1}$ map $\Z^2$ into $\Z^2$, so $\Phi$ restricts to a bijection
$P \cap \Z^2 \to \Phi(P) \cap \Z^2$ for every $P \subseteq \R^2$: the
lattice-point count is invariant under unimodular equivalence. But
$\bigl|[0,m]^2 \cap \Z^2\bigr| = (m+1)^2$ is a perfect square, and
$2^2 = 4 < 7 < 9 = 3^2$, so no $m$ gives $7$. With (ii), $\conv(S_A)$ is
equivalent to no $[0,m]^2$.

Finally, (i) and (iii) contradict the ``only if'' direction of the quoted
conjecture for $S = S_A$ --- under either reading of ``for any $m$'', since
(iii) denies the equivalence for every $m$ individually.
\end{proof}

\begin{remark}
$S_A$ is not the lattice-point set of its hull: it omits $(1,1)$, $(1,2)$,
$(2,3)$. Conjecture~3 of \cite{BDLT2025} is stated for arbitrary bounded
$S \subset \Z^d$, so this is no defect of the counterexample --- but it does
suggest the repair of restricting the conjecture to sets with
$S = \conv(S) \cap \Z^d$. Section~\ref{sec:repair} shows that this repair
also fails.
\end{remark}

\section{The full-lattice-set repair also fails}\label{sec:repair}

The results of this section are verified in exact integer and rational
arithmetic (Section~\ref{sec:verification}) but are \emph{not} part of the
Lean development.

\begin{theorem}[witness $T$, $d = 2$]\label{thm:witnessB}
Let $T = \{(0,1),\,(1,0),\,(1,1),\,(2,2)\} \subset \Z^2$. Then
$T = \conv(T) \cap \Z^2$, $\betaZ(T) = 4 = 2^2$, and $\conv(T)$ --- the
triangle with vertices $(0,1)$, $(1,0)$, $(2,2)$, of area $3/2$ --- is
unimodularly equivalent to no square $[0,m]^2$.
\end{theorem}

\begin{proof}
The six pairwise differences $(1,-1)$, $(1,0)$, $(2,1)$, $(0,1)$, $(1,2)$,
$(1,1)$ are primitive, so $\betaZ(T) = 4$ by Lemma~\ref{lem:beta}. Since
$(1,1) = \tfrac13\bigl[(0,1) + (1,0) + (2,2)\bigr]$ is the centroid of the
other three points, $\conv(T)$ is the triangle with the stated vertices; its
area is $3/2$ by the shoelace formula. The triangle is cut out by
$x + y \ge 1$, $2x - y \le 2$, $2y - x \le 2$ (each line through two
vertices, each inequality satisfied by all of $T$ and strictly by the
centroid $(1,1)$), and a four-line case check gives the integer solutions:
the second and third inequalities force $4x - 4 \le 2 + x$, i.e.\ $x \le 2$,
and $x \ge 0$ (for $x \le -1$ they give $2y \le 1$ and $y \ge 1 - x \ge 2$,
impossible); then $x = 0 \Rightarrow y = 1$, $x = 1 \Rightarrow y \in
\{0,1\}$, $x = 2 \Rightarrow y = 2$. So $\conv(T) \cap \Z^2 = T$, with $4$
points. If $\conv(T) = A\,[0,m]^2 + t$ were a unimodular image of a square,
the count invariant of Theorem~\ref{thm:witnessA}(iii) forces
$(m+1)^2 = 4$, i.e.\ $m = 1$; but $|{\det A}| = 1$, so unimodular maps
preserve area, and $\operatorname{area}\bigl(\conv(T)\bigr) = 3/2 \neq 1 =
\operatorname{area}\bigl([0,1]^2\bigr)$. (Alternatively: an affine bijection
preserves the number of vertices, and a triangle has $3 \neq 4$.)
\end{proof}

\begin{theorem}[witness $C$, $d = 3$]\label{thm:witnessC}
Let
\[
C = \{(0,2,2),\,(1,0,1),\,(1,1,2),\,(1,2,2),\,(2,1,1),\,(2,1,2),\,(2,2,1),
\,(3,1,3)\} \subset \Z^3 .
\]
Then $C = \conv(C) \cap \Z^3$, $\betaZ(C) = 8 = 2^3$, and $\conv(C)$, of
volume $5/2$, is unimodularly equivalent to no cube $[0,m]^3$.
\end{theorem}

\begin{proof}
All $28$ pairwise differences are primitive ($28$ exact gcd computations),
so $\betaZ(C) = 8$ by Lemma~\ref{lem:beta}. For $C = \conv(C) \cap \Z^3$ we
use an exact supporting-halfspace certificate: collect every plane spanned by
an affinely independent triple of points of $C$ having all eight points on
one (closed) side. Since $\conv(C)$ is full-dimensional (verified: four of
the points are affinely independent), every facet plane of $\conv(C)$ is
spanned by points of $C$ and hence occurs in this collection, so the
intersection of the collected halfspaces is exactly $\conv(C)$; enumerating
the integer points of the coordinate bounding box
$[0,3] \times [0,2] \times [1,3]$ (which contains $\conv(C)$) against these
inequalities returns precisely the eight points of $C$. The volume,
$5/2$, is computed exactly from the same facet description by fan
triangulation.

Now suppose $\conv(C) = A\,[0,m]^3 + t$ with $A \in \GL(3,\Z)$,
$t \in \Z^3$. The count invariant (as in Theorem~\ref{thm:witnessA}(iii),
verbatim in $d = 3$) gives $(m+1)^3 = |C| = 8$, so $m = 1$: indeed $m = 0$
gives $1$ point and $m \ge 2$ gives at least $27 > 8$. For $m = 1$ the map
$\Phi(x) = Ax + t$ restricts to a bijection from
$[0,1]^3 \cap \Z^3 = \{0,1\}^3$ onto $\conv(C) \cap \Z^3 = C$. For a finite
$X \subset \Z^3$ let
\[
N(X) \;=\; \#\bigl\{\, \{\{p,q\},\{r,s\}\} :
\{p,q\} \neq \{r,s\} \text{ pairs of distinct points of } X,\
p + q = r + s \,\bigr\},
\]
the number of unordered coincidences among the pairwise sums.
Since $A$ is invertible, $p + q = r + s$ iff
$\Phi(p) + \Phi(q) = \Phi(r) + \Phi(s)$, so $N$ is preserved by $\Phi$. But
$N(\{0,1\}^3) = 12$ --- the $\binom{4}{2} = 6$ coincidences among the four
antipodal vertex pairs, which all sum to $(1,1,1)$, plus one coincidence of
the two diagonals in each of the $6$ facets --- while $N(C) = 2$: the only
coinciding pair-sums in $C$ are
\begin{gather*}
(0,2,2) + (2,1,2) \;=\; (2,3,4) \;=\; (1,1,2) + (1,2,2),\\
(1,1,2) + (2,2,1) \;=\; (3,3,3) \;=\; (1,2,2) + (2,1,1).
\end{gather*}
This contradiction rules out $m = 1$. (Cross-check, independent of $N$:
unimodular maps preserve volume, and
$\vol\bigl(\conv(C)\bigr) = 5/2 \neq 1 = \vol\bigl([0,1]^3\bigr)$.)
\end{proof}

\begin{remark}[what survives]\label{rem:survives}
All three witnesses exploit the regime $\diamZ(S) = 1$, where
Lemma~\ref{lem:beta} collapses $\betaZ$ to the cardinality. We do not
determine the status of the conjecture under the combined repair
``$S = \conv(S) \cap \Z^d$ \emph{and} $\diamZ(S) \ge 2$'', which our
witnesses do not reach; for $d = 2$ the B\'ar\'any--F\"uredi
characterization \cite[Theorem~2(iii)]{BaranyFuredi2001} cited by the
authors is a natural route to a positive result there. We note only that the
needed diameter restriction is somewhat awkward for the conjecture's
motivation, since the sharpness example $\{0,1\}^d$ of
\cite[Theorem~4.8]{BDLT2025} itself has lattice diameter $1$. The ``if''
direction of the conjecture is not addressed by this note.
\end{remark}

\section{The Lean 4 artifact}\label{sec:lean}

Theorem~\ref{thm:witnessA} is formalized in Lean~4 \cite{Lean4} (toolchain
\texttt{leanprover/lean4:v4.30.0}) over mathlib \cite{mathlib} (release tag
\texttt{v4.30.0}), in a development of about $1100$ lines. The library
consists of seven files:

\begin{center}\small
\begin{tabular}{@{}lll@{}}
\toprule
File & Content & Lines\\
\midrule
\texttt{Borsuk/Defs.lean} & frozen definitions + faithfulness docstrings & 282\\
\texttt{Borsuk/SegmentGcd.lean} & Lemma~\ref{lem:gcd}: segment count $= \gcd + 1$ & 252\\
\texttt{Borsuk/WitnessA.lean} & Lemma~\ref{lem:beta}; $\diamZ(S_A) = 1$; $\betaZ(S_A) = 4$ & 148\\
\texttt{Borsuk/HullSeven.lean} & H-representation; $|\conv(S_A) \cap \Z^2| = 7$ & 176\\
\texttt{Borsuk/Unimodular.lean} & count invariance; $(m+1)^2$ cube count; $(m+1)^2 \neq 7$ & 167\\
\texttt{Borsuk/Main.lean} & assembly of the three main theorems & 59\\
\texttt{Borsuk/Smoke.lean} & toolchain smoke tests & 44\\
\bottomrule
\end{tabular}
\end{center}

\noindent The three top-level results, verbatim from
\texttt{Borsuk/Main.lean}:

\begin{lstlisting}
theorem witnessA_kills_conjecture3 :
    betaZ SA = 4 ∧
      ∀ m : ℕ, ¬ UnimodEquiv (convexHull ℝ (toR '' (SA : Set Z2))) (cube m)

theorem conjecture3_counterexample :
    ∃ S : Finset Z2, betaZ S = 2 ^ 2 ∧
      ∀ m : ℕ, ¬ UnimodEquiv (convexHull ℝ (toR '' (S : Set Z2))) (cube m)

theorem conjecture3_false : ¬ Conjecture3For2D
\end{lstlisting}

\noindent where the formal statement of (the $d = 2$ instance of) the
conjecture, from \texttt{Borsuk/Defs.lean}, is

\begin{lstlisting}
def Conjecture3For2D : Prop :=
  ∀ S : Finset Z2,
    (betaZ S = 2 ^ 2 ↔
      ∃ m : ℕ, UnimodEquiv (convexHull ℝ (toR '' (S : Set Z2))) (cube m))
\end{lstlisting}

\paragraph{Encoding choices (faithfulness).}
The Lean kernel verifies the proofs; whether the formal \emph{definitions}
say what the paper says is the one link a proof checker cannot certify, so
\texttt{Defs.lean} documents every definition against the verbatim \LaTeX{}
source of \cite{BDLT2025}, with the following choices. $\Z^2$ is
\texttt{Z2 = $\mathbb{Z}$ $\times$ $\mathbb{Z}$}, embedded in
$\R^2$ by the injection \texttt{toR}. A ``bounded set $S \subset \Z^d$'' is
a \texttt{Finset Z2} (bounded subsets of $\Z^2$ are exactly the finite
ones; the witness, a $4$-element set, is bounded under any reading).
$f$ is \texttt{latticeLength}, the cardinality of
$\{p \in \Z^2 : p \in [x,y]\}$ minus one, with mathlib's \texttt{segment}
proved equal to the paper's $\conv(\{x,y\})$;
$\diamZ$ is a double \texttt{Finset.sup}. Partitions are mathlib's
\texttt{Finpartition} (pairwise-disjoint nonempty parts with union $S$), and
\texttt{betaZ} is the infimum ($\mathtt{sInf}$) of the achievable part
counts of partitions all of whose parts have strictly smaller
\texttt{latticeDiam} --- the paper's ``minimal size of a partition \dots\
[with] strictly smaller lattice diameter''. Unimodular equivalence is an
explicit structure \texttt{UnimodMap}: four integer matrix entries with
$ad - bc = \pm 1$ (the units of $\Z$) plus an integer translation, acting on
$\R^2$, with $P = \varphi(Q)$ as the paper's $P = AQ + t$; the formal
development proves that such maps biject $\Z^2$ (the paper's ``lattice
preserving maps''). \texttt{cube m} is $[0,m]^2 \subseteq \R^2$, and
\texttt{latticeCount} is the number of lattice points. Two deliberate
robustness margins: ``for any $m \in \N$'' is rendered as
$\exists\, m$, but the proof establishes $\forall m,\ \neg\,\mathrm{equiv}$,
which negates the right-hand side under both the existential and the literal
universal reading, for $\N$ with or without $0$; and \texttt{UnimodEquiv} is
stated for arbitrary subsets of $\R^2$ rather than lattice polytopes only ---
on the objects in the conjecture ($\conv(S_A)$ and $[0,m]^2$, both lattice
polytopes) the readings coincide. One degenerate corner is flagged rather
than hidden: for the empty set the empty partition is admissible and gives
$\betaZ(\emptyset) = 0$, and for singleton $S$ no admissible partition
exists, so the achievable-count set is empty and the convention
$\mathtt{sInf}\,\emptyset = 0$ gives $\betaZ(S) = 0$ where the paper's
minimum is simply undefined; if moreover $0 \in \N$, a one-point hull
\emph{is} unimodularly equivalent to the degenerate cube $[0,0]^2$, so a
singleton by itself already falsifies the formal biconditional. The
development makes no use of this junk route: the proof of
\texttt{conjecture3\_false} instantiates the conjecture only at $S = S_A$,
with honest lattice diameter $1$ and an honestly attained minimum
$\betaZ(S_A) = 4$, and Theorem~\ref{thm:witnessA} refutes the ``only if''
direction in a form untouched by any convention for the degenerate cases.

\paragraph{Verification status.}
On 2026-06-11 and again on 2026-06-12 the project built cleanly:
\texttt{lake build} reports \texttt{Build completed successfully (8483
jobs)}, exit code $0$, no warnings. There are no \texttt{sorry} or
\texttt{admit} terms, no \texttt{axiom} declarations, no
\texttt{native\_decide} (all decidability checks are kernel-checked
\texttt{decide}), no \texttt{unsafe} or \texttt{partial} definitions, and no
custom elaboration. The axiom audit (\texttt{lake env lean
scripts/CheckAxioms.lean}) prints, for all sixteen audited theorems ---
including the three above and the supporting lemmas
\texttt{ncard\_segLatticePts}, \texttt{latticeDiam\_SA},
\texttt{betaZ\_eq\_card\_of\_latticeDiam\_eq\_one}, \texttt{betaZ\_SA},
\texttt{mem\_hullSA\_iff}, \texttt{latticeCount\_hullSA},
\texttt{UnimodEquiv.latticeCount\_eq}, \texttt{latticeCount\_cube},
\texttt{succ\_sq\_ne\_seven} --- an axiom set contained in
\[
[\texttt{propext},\ \texttt{Classical.choice},\ \texttt{Quot.sound}],
\]
the three standard mathlib axioms; in particular no \texttt{sorryAx}. For
example (one line per theorem, quoted exactly):
{\scriptsize
\begin{verbatim}
'Borsuk.witnessA_kills_conjecture3' depends on axioms: [propext, Classical.choice, Quot.sound]
'Borsuk.conjecture3_false' depends on axioms: [propext, Classical.choice, Quot.sound]
\end{verbatim}
}

\paragraph{Building.}
With elan installed, from the \texttt{lean/} directory of the artifact:
\begin{quote}\small
\begin{verbatim}
elan toolchain install leanprover/lean4:v4.30.0
lake exe cache get      # fetch mathlib v4.30.0 build cache (~1.6 GB)
lake build              # expect: Build completed successfully, exit 0
lake env lean scripts/CheckAxioms.lean   # axiom audit, output as above
\end{verbatim}
\end{quote}
The toolchain and mathlib revisions are pinned by \texttt{lean-toolchain},
\texttt{lakefile.toml}, and \texttt{lake-manifest.json}, so the build is
reproducible; without the cache, mathlib builds from source in a few hours.

\section{Verification of the numbers in this note}\label{sec:verification}

Every quantity quoted above is mechanically checkable. For
Theorem~\ref{thm:witnessA} the authority is the Lean development of
Section~\ref{sec:lean}. For Theorems~\ref{thm:witnessB}
and~\ref{thm:witnessC}, and as an independent cross-check of the witness-$A$
data, the script \texttt{\_recompute.py} distributed with this note (Python
standard library only; integer and \texttt{fractions.Fraction} arithmetic on
every accept path, no floating point anywhere) re-verifies by direct
recomputation: the $6 + 6 + 28$ primitivity checks; the
H-representation of $\conv(S_A)$ (halfspace containment and the edge--vertex
incidences), the enumeration of its $7$ lattice points, and exact
convex-combination certificates for the three non-elements of $S_A$; the
triangle $\conv(T)$, its area $3/2$, and $\conv(T) \cap \Z^2 = T$; the
supporting-halfspace certificate for $\conv(C) \cap \Z^3 = C$ (these
halfspaces the script does derive from scratch), the exact
volume $5/2$, the pair-sum invariants $N(C) = 2$ and $N(\{0,1\}^3) = 12$,
and the cube counts $(m+1)^3$. It was re-run on 2026-06-12 with verdict
\texttt{ALL CHECKS PASS} ($27$ checks). A second script,
\texttt{\_referee\_recheck.py}, written during the hostile-referee pass of
2026-06-12 deliberately by different routes --- hull membership by exact
Carath\'eodory solves over $\mathbb{Q}$ instead of H-representations,
vertices by the ``not in the hull of the others'' criterion, the volume by a
divergence-style signed surface sum, and a full census of coinciding
pair-sums --- reaches the same conclusions ($37$ checks,
\texttt{ALL CHECKS PASS}). The discovery pipeline had already verified
witnesses $A$--$C$ in exact arithmetic (three independent recomputations,
plus the second-model-family re-verification noted in
Section~\ref{sec:intro}); both scripts above were written separately from
it.

\paragraph{Data availability.}
The Lean project (source, pinned manifests, axiom-audit script), the two
exact-arithmetic verification scripts, the discovery-time verification
records, and the frozen copy of the arXiv v1 source of \cite{BDLT2025}
against which the definitions were checked are available at
\url{https://doi.org/10.5281/zenodo.20673865} and
\url{https://github.com/demonstrandum-research/artifacts}.

\begin{thebibliography}{9}

\bibitem{BaranyFuredi2001}
I.~B\'ar\'any and Z.~F\"uredi,
\emph{On the lattice diameter of a convex polygon},
Discrete Math.\ \textbf{241} (2001), 41--50.

\bibitem{Borsuk1933}
K.~Borsuk,
\emph{Drei S\"atze \"uber die $n$-dimensionale euklidische Sph\"are},
Fund.\ Math.\ \textbf{20} (1933), 177--190.

\bibitem{BDLT2025}
A.~E. Brose, J.~A. De Loera, G.~Lopez-Campos, and A.~J. Torres,
\emph{On lattice diameter segments and a discrete Borsuk partition problem},
arXiv:2508.20009 (v1, 27 August 2025).

\bibitem{KahnKalai1993}
J.~Kahn and G.~Kalai,
\emph{A counterexample to Borsuk's conjecture},
Bull.\ Amer.\ Math.\ Soc.\ \textbf{29} (1993), 60--62.

\bibitem{mathlib}
The mathlib Community,
\emph{The Lean mathematical library},
in: Proceedings of the 9th ACM SIGPLAN International Conference on Certified
Programs and Proofs (CPP 2020), ACM, 2020, pp.~367--381.

\bibitem{Lean4}
L.~de Moura and S.~Ullrich,
\emph{The Lean 4 theorem prover and programming language},
in: Automated Deduction -- CADE 28, Lecture Notes in Comput.\ Sci.\
\textbf{12699}, Springer, 2021, pp.~625--635.

\bibitem{Rabinowitz1989}
S.~Rabinowitz,
\emph{A theorem about collinear lattice points},
Utilitas Math.\ \textbf{36} (1989), 93--95.

\end{thebibliography}

\end{document}
